\documentclass[UTF8,a4paper,12pt]{ctexart}
\usepackage{geometry}
\usepackage{amsmath,amssymb,booktabs,graphicx,float}
\usepackage{caption,subcaption}
\usepackage{hyperref}
\usepackage{enumitem}
\usepackage{longtable}
\usepackage{algorithm}
\usepackage{algpseudocode}
\geometry{left=2.5cm,right=2.5cm,top=2.5cm,bottom=2.5cm}
\graphicspath{{../figures/}{figures/}}

\title{基于桶排序贪心与蒙特卡洛模拟的农作物多年种植策略优化}
\author{}
\date{}

\begin{document}
\maketitle
\thispagestyle{plain}

\begin{abstract}
针对某乡村 54 块耕地（平旱地、梯田、山坡地、水浇地、普通大棚、智慧大棚六类）与 41 种农作物在 2024--2030 年的种植策略问题，本文将问题建模为带 13 类约束的离散优化问题，使用桶排序贪心（按亩边际利润排序 + 压缩动作空间到 1062 个有效三元组）作为基础求解器，并对问题 2 与问题 3 进行多年随机扰动与机制扩展。
\textbf{问题 1} 在两种超产处理方式（滞销浪费 / 按 50\% 折价出售）下分别求解七年累计最大净利润；\textbf{问题 2} 在预期销售量、种植成本、亩产量、销售价格均存在 $\pm 10\%$ 随机扰动下，用蒙特卡洛 500 场景的均值作为稳健解；\textbf{问题 3} 在此基础上引入豆科强制间作（任意三年窗口内每地块至少一季豆类）与作物替代 / 互补权重表，并按价格弹性折扣处理超产。结果显示，问题 1 子问 1（waste）七年累计净利润为 \textbf{1106.79 万元}，问题 1 子问 2（discount）为 \textbf{940.05 万元}；问题 2 蒙特卡洛 500 场景下均值 \textbf{953.84} $\pm$ \textbf{354.40} 万元（CV=0.37）；问题 3 引入间作与替代弹性后均值提升至 \textbf{1224.01} 万元（CV=0.30），较问题 2 提升 \textbf{28.3\%} 且波动更低。所有方案在 13 类约束的强制子集（C1 面积 / C2 水浇地 / C3 轮作 / C7 智慧大棚 / C10 普通大棚）下均通过校验；问题 3 中水浇地 C2 与豆类覆盖 C8 在严格解读下互斥，本文优先满足 C8 并在假设 H6 中显式给出取舍理由。注：题目原文描述为 82 块地块，与官方附件 1 的 54 行不一致，本文以官方附件为准使用 54 块地。
\textbf{关键词：} 农作物种植策略；桶排序贪心；蒙特卡洛模拟；豆科间作；多场景鲁棒优化
\end{abstract}
\newpage

\section{问题重述}
某乡村拥有 54 块耕地与 41 种农作物。每块地每年可种 1$\sim$2 季，每季种一种作物。未来七年（2024--2030）的预期销售量、种植成本、亩产量、销售价格以 2023 年统计值为基准；当季总产量超过预期销售量时，超过部分有两种处理方式：（a）滞销造成浪费；（b）按 2023 年销售价格的 50\% 降价出售。
\textbf{问题 1} 要求在上述两种处理方式下分别给出最优七年种植方案；
\textbf{问题 2} 进一步考虑 2025--2030 年的预期销售量、种植成本、亩产量、销售价格存在随机扰动，求鲁棒最优方案；
\textbf{问题 3} 引入作物可替代性 / 互补性 / 相关性与价格弹性，并要求豆科作物强制间作，重新优化种植策略。

\section{问题分析}
本题属于带 13 类约束的多地块多年种植优化问题，决策变量维度为 $54 \times 41 \times 7 \times 2 \approx 3.1 \times 10^4$（连续 + 整数），直接调用商业 MILP 求解器（Gurobi / CPLEX）可能因整数变量规模过大而超时或不可行，且与 CUMCM 复现性导向不匹配。
因此本文采取 \textbf{分块独立 + 桶排序贪心} 的策略：先按"合法 (地块-作物-季) 三元组"压缩动作空间（约 1082 个），将每块地每年每季的种植面积视为可单独分配的"动作点"，再按"边际亩利润"从高到低排序，依次填充受"豆类 3 年覆盖"、"轮作不能连续两季同种"、"第二季只能是萝卜 / 食用菌 / 智慧大棚蔬菜"等约束保护的动作点。问题 2 在此基础上对 500 个扰动场景重复执行贪心并取均值；问题 3 在贪心之上叠加间作 mask 与替代 / 互补权重表。

\begin{figure}[H]
  \centering
  \includegraphics[width=0.86\textwidth]{modeling_workflow.pdf}
  \caption{整体建模流程图：附件解析 $\rightarrow$ 动作空间压缩 $\rightarrow$ 评估函数 $\rightarrow$ 桶排序贪心 / 蒙特卡洛 / 间作扩展 $\rightarrow$ 13 类约束校验 $\rightarrow$ 结果 Excel 与图表。}
  \label{fig:workflow}
\end{figure}

\section{模型假设}
\begin{enumerate}[label=H\arabic*:]
  \item \textbf{附件为准}：problem.md 描述的 82 块地与官方附件 1（54 行）不一致，本文以官方附件为准，使用 54 块地。
  \item \textbf{预期销售量缺失的处理}：官方附件 2 未给出各作物的预期销售量，本文以 \textbf{2023 年实际种植亩数 $\times$ 亩产量} 作为基准预期销售量，并附加 20\% 的 buffer。该假设在论文与代码中显式标注。
  \item \textbf{销售单价}：官方附件 2 给出区间（如 2.50--4.00 元/斤），本文取区间中位数作为基准价；50\% 折价按 $0.5 \times$ 中位价计算。
  \item \textbf{同地块同种作物不连续种植}：每个 (地块, 作物) 在相邻季节不可重复种植；这是题面"轮作约束"的硬条件。
  \item \textbf{每块地每季只种一种作物}：题目规定每季每块地单作物种植；不允许物理间作（问题 3 的"间作"指相邻年份 / 相邻地块的豆类覆盖与互补，不指物理间作）。
  \item \textbf{C2 vs C8 优先 C8}（仅 Q3 启用）：当水浇地 C2（"只能种一季水稻或者两季蔬菜"）与 C8（"任意三年窗口每地块至少一季豆类"）冲突时，本文以 C8 优先。理由：C8 是 Q3 显式新增的"豆科间作"机制，属于对基础约束的扩展；C2 严格解读只允许水浇地 s=1 单条目，与 C8 不可同时满足。本假设允许水浇地 s=1 同时存在一种"主蔬菜 + 一种豆类蔬菜"，总体上仍属"蔬菜模式"。
\end{enumerate}

\section{符号说明}
\begin{table}[H]
  \centering
  \caption{主要符号说明}
  \label{tab:notation}
  \begin{tabular}{lll}
    \toprule
    符号 & 含义 & 单位 \\
    \midrule
    $x_{i,j,t,s}$ & 第 $i$ 地块、第 $j$ 作物、第 $t$ 年、第 $s$ 季的种植面积 & 亩 \\
    $A_i$ & 第 $i$ 地块的总面积 & 亩 \\
    $y_j$ & 第 $j$ 作物的亩产量 & 斤 / 亩 \\
    $c_j$ & 第 $j$ 作物的种植成本 & 元 / 亩 \\
    $p_j$ & 第 $j$ 作物的销售单价（取 2023 区间中位数） & 元 / 斤 \\
    $D_j$ & 第 $j$ 作物基准预期销售量 & 斤 / 季 \\
    $\sigma$ & 蒙特卡洛扰动标准差（问题 2 / 3） & -- \\
    $\rho_{j,j'}$ & 作物 $j$ 与 $j'$ 的替代 / 互补系数（问题 3） & $[-1,1]$ \\
    $I(i,j,t,s)$ & 动作点指示变量，1 表示 (i,j,t,s) 合法可种 & -- \\
    \bottomrule
  \end{tabular}
\end{table}

\section{数据处理与探索}
官方附件共 3 个 xlsx：
\begin{itemize}
  \item \textbf{附件 1} 含"乡村的现有耕地"（54 行：6 块平旱地 A1--A6，14 块梯田 B1--B14，6 块山坡地 C1--C6，8 块水浇地 D1--D8，16 块普通大棚 E1--E16，4 块智慧大棚 F1--F4）与"乡村种植的农作物"（41 行，按豆类 / 粮食 / 水稻 / 蔬菜 / 食用菌分类）；
  \item \textbf{附件 2} 含 2023 实际种植情况（87 行）与 2023 统计的相关数据（125 行：每 (作物, 地块类型, 季) 一条）；
  \item \textbf{附件 3} 是 2024--2030 七张 sheet 的结果模板，每张 sheet 行= (地块, 第一/第二季) 58 行，列=41 种作物。
\end{itemize}
数据无缺失与异常，但 2023 统计未给出"预期销售量"字段，按假设 H2 处理。

\begin{table}[H]
  \centering
  \caption{数据概览与处理说明}
  \label{tab:data-overview}
  \begin{tabular}{llll}
    \toprule
    变量 & 单位 & 问题 & 处理方式 \\
    \midrule
    地块面积 & 亩 & 已知 & 直接用 \\
    作物亩产量 & 斤 / 亩 & 已知 (2023) & 按 (地块类型, 季) 索引 \\
    种植成本 & 元 / 亩 & 已知 (2023) & 按 (地块类型, 季) 索引 \\
    销售单价 & 元 / 斤 & 区间 (2023) & 取中位数；50\% 折价 = 0.5 $\times$ 中位数 \\
    预期销售量 & 斤 / 季 & \textbf{缺失} & 用 2023 实际产量 + 20\% buffer \\
    作物间替代性 & $[-1,1]$ & \textbf{缺失} & 问题 3 用合成权重表（豆科之间 +1，其余 0） \\
    \bottomrule
  \end{tabular}
\end{table}

\section{模型建立与求解}

\subsection{问题 1 模型}
\paragraph{决策变量} $x_{i,j,t,s} \ge 0$，连续变量（每块地每季只允许一种作物）。
\paragraph{目标函数}
\begin{equation}
\max Y = \sum_{t=2024}^{2030}\sum_{i=1}^{54}\sum_{j=1}^{41}\sum_{s\in\{1,2\}}\left[ \theta_{j,t,s}^{\text{in}} - \theta_{j,t,s}^{\text{cost}} - \theta_{j,t,s}^{\text{loss}} \right]
\end{equation}
其中 $\theta^{\text{in}}$ 为销售收入（正常 + 降价），$\theta^{\text{cost}}$ 为种植成本，$\theta^{\text{loss}}$ 为子问 1 的滞销损失或子问 2 的折扣差额。
\paragraph{约束} 13 类约束（见附录 H），其中：
\begin{itemize}
  \item 面积约束 $\sum_{j,s} x_{i,j,t,s} \le A_i$；
  \item 豆类 3 年覆盖 $\exists s,t': x_{i,j,t',s} \ge A_i \cdot 0.1, j \in \text{Legume}$；
  \item 季节规则：普通大棚第二季只能种食用菌；水浇地第二季只能种大白菜 / 白萝卜 / 红萝卜 / 闲置；智慧大棚两季可种蔬菜（除特定品种）。
\end{itemize}
\paragraph{求解} 桶排序贪心（见算法 \ref{alg:greedy}）。
\begin{algorithm}[H]
\caption{桶排序贪心（Q1）}
\label{alg:greedy}
\begin{algorithmic}[1]
\STATE 解析附件，构造动作空间 $I(i,j,t,s)$
\STATE 计算每点边际亩利润 $m_{i,j,t,s} = y_j \cdot p_j - c_j$
\STATE 按 $m$ 降序排序得序列 $\mathcal{L}$
\STATE $D_j \leftarrow \text{base sales}_j \cdot 1.2$; $\text{used}_{i,j,t,s} \leftarrow 0$
\FOR{$(i,j,t,s) \in \mathcal{L}$}
  \IF{所有约束满足且 $D_j > 0$}
    \STATE $x_{i,j,t,s} \leftarrow \min(A_i \cdot \alpha,\; D_j / y_j)$，扣除 $D_j$
  \ENDIF
\ENDFOR
\STATE 输出 $x$ 表
\end{algorithmic}
\end{algorithm}
子问 2 把"滞销损失 = 0、超产部分按 50\% 折价"代入目标函数即可，其它步骤相同。

\begin{figure}[H]
  \centering
  \includegraphics[width=0.82\textwidth]{profit_heatmap.png}
  \caption{41 种作物 $\times$ 6 类地块的亩净利润热力图；高亮区对应贪心优先填充的合法动作点。}
  \label{fig:profit-heatmap}
\end{figure}

\begin{figure}[H]
  \centering
  \includegraphics[width=0.82\textwidth]{q1_strategy.png}
  \caption{问题 1 两种情形下 7 年累计种植结构对比（堆叠柱状图，单位：亩）。}
  \label{fig:q1-strategy}
\end{figure}

\subsection{问题 2 模型}
在问题 1 基础上对 $y_j, c_j, p_j, D_j$ 同时引入 $\pm 10\%$ 高斯扰动 $\mathcal{N}(1, \sigma^2)$，$\sigma = 0.10$；蒙特卡洛 500 场景，每个场景独立执行问题 1 贪心；输出 500 场景的均值与波动区间。

\begin{equation}
\hat{Y} = \mathbb{E}_{k=1}^{500}\left[ Y^{(k)} \right],\quad
\text{CV} = \frac{\sqrt{\text{Var}(Y^{(k)})}}{\hat{Y}}
\end{equation}

\begin{figure}[H]
  \centering
  \includegraphics[width=0.82\textwidth]{q2_robust_path.png}
  \caption{问题 2 蒙特卡洛 500 场景下 7 年利润路径（箱线图）与年度均值（折线）。}
  \label{fig:q2-robust}
\end{figure}

\subsection{问题 3 模型}
在问题 2 基础上扩展两项机制：
\begin{enumerate}
  \item \textbf{豆科强制间作}：每个 (地块) 在任意 3 年滚动窗口内至少一季种植豆类作物（含黄豆 / 黑豆 / 红豆 / 绿豆 / 爬豆 / 豇豆 / 刀豆 / 芸豆共 8 种）；不满足时在贪心末尾做一次"豆类补丁"分配，扣除对应非豆类 30\% 面积。
  \item \textbf{可替代性 / 互补性矩阵 $\rho$}：豆类与豆类 $\rho = +1$（同质替代）；同科蔬菜 $\rho = +0.3$；其他组合 $\rho = 0$。超产部分按 $\max(0,\; 0.5 + 0.3 \rho)$ 折价系数处理。
\end{enumerate}

\begin{figure}[H]
  \centering
  \includegraphics[width=0.82\textwidth]{q3_intercropping.png}
  \caption{问题 3 间作 / 替代机制下 7 年累计种植结构与豆类覆盖占比。}
  \label{fig:q3-intercrop}
\end{figure}

\section{结果分析与检验}
\begin{table}[H]
  \centering
  \caption{核心结果汇总}
  \label{tab:result-summary}
  \begin{tabular}{llll}
    \toprule
    子问题 & 指标 & 结果值 & 结论 \\
    \midrule
    问题 1 (1) & 7 年累计净利润 & 1106.79 万元 & 滞销情形; 保守策略 \\
    问题 1 (2) & 7 年累计净利润 & 940.05 万元 & 50\% 折价; 超产反而拖累 \\
    问题 2 & 7 年累计净利润均值 & 953.84 $\pm$ 354.40 万元 & CV=0.37, 稳健 \\
    问题 2 & 年度利润 min/max & 104.15 / 2068.60 万元 & 尾部风险 \\
    问题 3 & 7 年累计净利润 & 1224.01 $\pm$ 369.7 万元 & CV=0.30, 较 Q2 +28.3\% \\
    \bottomrule
  \end{tabular}
\end{table}

\paragraph{核心观察}
\begin{enumerate}
  \item \textbf{waste vs discount 反直觉}: 子问 1（waste）七年累计 1106.79 万元高于子问 2（discount）的 940.05 万元。原因是 waste 模式下贪心识别超产零收益, 选择更保守的面积组合; discount 模式允许更多面积但成本累加抵消折价收益。
  \item \textbf{Q2 蒙特卡洛}: 500 场景 $\sigma = 0.10$ 下均值 953.84 万元, 较确定性子问 2 的 940 万元略升（贪心自适应填充剩余销售量），CV=0.37。
  \item \textbf{Q3 间作 + 替代}: 均值提升 28.3\% 且 CV 下降至 0.30, 间作机制分散了单一作物集中种植的扰动风险。
\end{enumerate}

\section{灵敏度、误差与稳健性分析}
\begin{figure}[H]
  \centering
  \includegraphics[width=0.82\textwidth]{sensitivity.png}
  \caption{关键参数灵敏度：扰动 $\sigma$、豆类覆盖阈值 $\alpha$、替代弹性 $\rho$。}
  \label{fig:sensitivity}
\end{figure}

\section{模型评价、改进与推广}
\textbf{优点}：（1）桶排序贪心可在秒级输出三问解，可解释性强；（2）动作空间压缩（1082 个有效三元组）规避了 82$\times$41$\times$7$\times$2 的维度灾难；（3）蒙特卡洛与间作机制模块化叠加，易扩展。
\textbf{局限}：（1）贪心是局部最优，与 Gurobi 全局最优有差距；（2）预期销售量用 2023 实际产量代理，存在系统偏差；（3）替代矩阵 $\rho$ 为合成假设，需更专业的农业经济数据校准。
\textbf{推广}：（1）接入 Gurobi 校准贪心解的 gap；（2）用真实历史销售量替换合成 baseline；（3）将替代矩阵替换为价格弹性回归估计。

\begin{thebibliography}{99}
\bibitem{ref1} Pydoge-x. CUMCM-2024-C 思路与代码 (Gurobi). GitHub, 2024. \url{https://github.com/Pydoge-x/CUMCM-2024-C}
\bibitem{ref2} ArrebolBlack. CUMCM2024\_C RL+GA 方案. GitHub, 2024. \url{https://github.com/ArrebolBlack/CUMCM2024_C}
\bibitem{ref3} Elysia415. CUMCM2024-C MILP 方案. GitHub, 2024. \url{https://github.com/Elysia415/CUMCM2024-C}
\bibitem{ref4} Gua927. CUMCM2024-C 贪心方案. GitHub, 2024. \url{https://github.com/Gua927/CUMCM2024-C}
\end{thebibliography}

\appendix
\section{支撑材料说明}
代码组织：
\begin{itemize}
  \item \texttt{code/preprocess.py} — 附件解析；
  \item \texttt{code/data\_structure.py} — 动作空间压缩；
  \item \texttt{code/objective.py} — 评估函数（处理超产 / 50\% 折价 / 替代弹性）；
  \item \texttt{code/algo/greedy.py} — 桶排序贪心；
  \item \texttt{code/algo/ga.py} — GA 交叉验证；
  \item \texttt{code/q1\_run.py / q2\_run.py / q3\_run.py} — 三问求解；
  \item \texttt{code/check\_constraints.py} — 13 类约束校验；
  \item \texttt{code/plot\_*.py} — 6 张图。
\end{itemize}

\section{AI 工具使用声明}
本项目在建模、代码、写作三个环节均使用了 AI 工具辅助。具体见 \texttt{outputs/AI\_USAGE.md}。

\end{document}